% Texas Tech Dissertation LaTeX Template 
% Adapted from University of Michigan Dissertation LaTeX Template, by John Meluso and Shreyas Kousik
% Modifications to follow TTU template by Huyen Nguyen and Vung Pham
% June 2023

\documentclass[thesis]{thesis-ttu}
%%% Packages included in thesis-umich class
%\RequirePackage[margin=1in,footskip=8pt,headsep=0.4cm,headheight=\baselineskip]{geometry}
%\RequirePackage{amsmath}
%\RequirePackage{amsfonts}
%\RequirePackage{amssymb}
%\RequirePackage{graphicx}
%\RequirePackage{subcaption}
%\RequirePackage{times}
% \RequirePackage{natbib}
%\RequirePackage{verbatim}
%\RequirePackage{upquote}
%\RequirePackage{textcomp}
%\RequirePackage{setspace}
%\RequirePackage{ifthen}
%\RequirePackage{soul}
%\RequirePackage{float}
%\RequirePackage[printonlyused]{acronym}
%\RequirePackage{makeidx}
\RequirePackage{fancyhdr}
%\RequirePackage{multicol}


% Include your packages here by including a \usepackage{<package_name>} command.
\usepackage{blindtext} % Example package which populates the ever-popular lorem ipsum text.
\usepackage[T1]{fontenc}
\usepackage{enumitem}
\usepackage{tabularx}
\usepackage{multirow}
\usepackage{url}
\usepackage{algorithm}
\usepackage[noend]{algpseudocode}
\usepackage{hyperref}
\usepackage{bibentry}      %  full citation in the body of the text (turn off natbib if use it)
\usepackage{graphicx}
\usepackage{pdflscape}
\usepackage{tikz}
\usepackage{rotating}
\usepackage{changepage}
\usepackage{geometry}

\usepackage{titlesec}

\titleformat{\section}
  {\fontfamily{phv}\selectfont \fontsize{13}{13}\bfseries}{\thesection}{1em}{}
\titleformat{\subsection}
  {\fontfamily{ptm}\selectfont \fontsize{13}{13}\bfseries}{\thesubsection}{1em}{} % import all packages from packages.tex

% Keep these next three lines in your preamble to use the alpha bibliography style, so that the references are left-aligned
\makeatletter
\renewcommand{\@biblabel}[1]{[#1]\hfill}
\makeatother

% Title of the dissertation
    \title{<Title of work, centered. Title \textbf{may not exceed} 238 characters (including spaces). \textcolor{blue}{Use consistent capitalization (first word only or all nouns). Single space all titles longer than one line. Do not bold.}>}

% Author name
\author{<First Name Last Name>}
\newcommand{\currentdegree}{<Initials for Currently Held Degrees>}

% Department
\department{<Name of Major>\\
\textcolor{blue}{\{official major listed on your degree\}}}

% Degree
\degree{DEGREE\\\textcolor{blue}{\{MASTER OF ARTS, MASTER OF SCIENCE, DOCTOR OF PHILOSOPHY, etc.\}}}

% Year of completion
\renewcommand{\year}{<Year of Graduation>}
\newcommand{\Month}{<Month>}

\rhead{Overleaf}

% Default style for front pages
\frontpagestyle{7} 

% Dedication
\hidededication
% \dedication[7]{\input{front_materials/dedication}}

% Acknowledgments 
\acknowledgments[7]{The purpose of this page is to recognize scholarly and professional aid and advice; however, the inclusion of references to persons who provided clerical help, help with field studies, financial assistance, and permission to use copyrighted materials is also acceptable.

Acknowledgments should be brief, in a professional style, and should not exceed two pages. 

\textit{\textcolor{blue}{\{The first page after your copyright page should be numbered ‘ii’. Your running header should begin on page ‘ii’.\}}}



}

% Preface
% \preface[7]{\input{front_materials/preface}}

% Committee
\committee{ %
<Name of Chair/Co-Chair, Degree initials> \\
Chair of the Committee\\

\vfill%
<Name of Committee Member/Co-Chair, Degree initials> \\

\vfill%
<Name of Committee Member, Degree initials> \\

\vfill%
<Name of Committee Member, Degree initials> \\


\vfill%
Mark Sheridan, Ph.D. \\
Dean of the Graduate School

}
% Commands to hide or show lists of figures, tables, etc.
\showlistoftables
% \showlistofprograms
% \showlistofappendices
\showlistofacronyms


% Definition of any acronyms used.
% To add an acronym, add an \acro{}{} command on a new line within the \acronyms{} command. For \acro, field 1 is the acronym and field 2 is the corresponding expression. See example below.

\acronyms{
    \acro{DOM}{Document Object Model}
    \acro{NLP}{Natural Language Processing}
    \acro{POS}{Part-of-speech}
}


% Some abstract text
\abstract{
Typically, an abstract will NOT have citations and should be no longer than 2-3 paragraphs.

\textit{\textcolor{blue}{\{The Abstract should follow the Table of Contents in the order of your front matter.\}}}

}
%\hideabstractpagenumber


%% DOCUMENT AREA
\begin{document}

\chapter{INTRODUCTION (HEADING 1)}
\label{ch:introduction}

The following is filler text (from “Pride and Prejudice” Jane Austen with apologies).

Elizabeth's mind was too full for conversation, but she saw and admired every remarkable spot and point of view. They gradually ascended for half-a-mile, and then found themselves at the top of a considerable eminence, where the wood ceased, and the eye was instantly caught by Pemberley House, situated on the opposite side of a valley, into which the road with some abruptness wound. It was a large, handsome stone building, standing well on rising ground, and backed by a ridge of high woody hills; and in front, a stream of some natural importance was swelled into greater, but without any artificial appearance. Its banks were neither formal nor falsely adorned. Elizabeth was delighted. She had never seen a place for which nature had done more, or where natural beauty had been so little counteracted by an awkward taste. They were all of them warm in their admiration; and at that moment she felt that to be mistress of Pemberley might be something!

\section{Examples}

\subsection{Figures}

You can include figures with image types PNG, JPEG etc. or PDF format. See Figures~\ref{fig:TTU} and Figure~\ref{fig:DoubleT}.

\begin{figure}[!htb]
    \centering
    \includegraphics[width=\linewidth]{chapters/figures/TTU.png}
    \caption{You can have a figure in PNG format like this one. \textcolor{red}{All figure titles that are longer than one line should be single-spaced and left aligned.  Figure titles that are one line only can be centered. Figures should be numbered and titled. The figure title may appear ABOVE OR BELOW the figure, but the placement should be consistent throughout your document.}}
    \label{fig:TTU}    
\end{figure}

\begin{figure}[!htb]
    \centering
    \includegraphics[width=0.5\linewidth]{chapters/figures/DoubleT.pdf}
    \caption{This is a figure in PDF format, it helps with high-resolution rendering. You can change the width of figure too, this one is at 0.5 line width.}
    \label{fig:DoubleT}    
\end{figure}

\subsection{Tables}

Please see \href{https://www.overleaf.com/learn/latex/Tables#Creating_a_simple_table_in_LaTeX}{overleaf.com/learn/latex/Tables} for more examples on tables. They descended the hill, crossed the bridge, and drove to the door; and, while examining the nearer aspect of the house, all her apprehension of meeting its owner returned. She dreaded lest the chambermaid had been mistaken. See Table~\ref{tab:my-table1}.

\textcolor{blue}{Tables should be numbered and titled. The table title should appear ABOVE the table.} 

\begin{table}[!htb]
    \caption{Country List. \textcolor{red}{All table titles that are longer than one line should be single-spaced and left aligned. Table titles that are only one line can be centered.}}
    \label{tab:my-table1}
\begin{tabular}{ |p{3cm}||p{3cm}|p{3cm}|p{3cm}|  }
 \hline
 \multicolumn{4}{|c|}{Country List} \\
 \hline
 Country Name or Area Name& ISO ALPHA 2 Code &ISO ALPHA 3 Code&ISO numeric Code\\
 \hline
 Afghanistan   & AF    &AFG&   004\\
 Aland Islands&   AX  & ALA   &248\\
 Albania &AL & ALB&  008\\
 Algeria    &DZ & DZA&  012\\
 American Samoa&   AS  & ASM&016\\
 \hline
\end{tabular}
\end{table}


\begin{table}[!htb]
\centering
\caption{Table titles that are only one line can be centered.}
\label{tab:my-table2}
\begin{tabular}{r r r r} 
 \hline
 Col1 & Col2 & Col2 & Col3 \\ [0.5ex] 
 \hline\hline
 1 & 6 & 87837 & 787 \\ 
 2 & 7 & 78 & 5415 \\
 3 & 545 & 778 & 7507 \\
 4 & 545 & 18744 & 7560 \\
 5 & 88 & 788 & 6344 \\ [1ex] 
 \hline
\end{tabular}

\end{table}

\subsection{Acronyms}

Here is an example of acronyms. The first time use is like this: \ac{DOM}. From the second time, it is: \ac{DOM}. This entry will then be included in the List of Acronyms page.

\section{Second-level subheading}


Mrs. Reynolds anticipated Miss Darcy's delight, when she should enter the room. "And this is always the way with him," she added. "Whatever can give his sister any pleasure is sure to be done in a moment. There is nothing he would not do for her."


\subsection{Third-level subheading}

The picture-gallery, and two or three of the principal bedrooms, were all that remained to be shown. In the former were many good paintings; but Elizabeth knew nothing of the art; and from such as had been already visible below, she had willingly turned to look at some drawings of Miss Darcy's, in crayons, whose subjects were usually more interesting, and also more intelligible.

In the gallery there were many family portraits, but they could have little to fix the attention of a stranger. Elizabeth walked in quest of the only face whose features would be known to her. At last it arrested her—and she beheld a striking resemblance to Mr. Darcy, with such a smile over the face as she remembered to have sometimes seen when he looked at her. She stood several minutes before the picture, in earnest contemplation, and returned to it again before they quitted the gallery. Mrs. Reynolds informed them that it had been taken in his father's lifetime.

\textcolor{blue}{Use a landscape-oriented page layout for wider tables.} See Table~\ref{tab:my-table3}.

\begin{landscape}
\centering
\thispagestyle{rotate}
\begin{table}[!htb]
\centering
\caption{Large, landscape-oriented table picked randomly from US Census. }
\label{tab:my-table3}
\begin{tabular}{|l|l|r|r|r|}

\hline
\textbf{NAICS code} & \textbf{Description}                                          & \textbf{Establishments} & \textbf{Revenue (\$1,000)} & \textbf{Paid employees} \\ \hline
524    & Insurance carriers \& related activities             & 9,412 & N         & 171,557     \\ \hline
5241   & Insurance   carriers                                 & 1,734 & Q         & 107,733     \\ \hline
52411  & Direct life, health, \& medical   insurance carriers & 735   & Q         & 66,152      \\ \hline
524113 & Direct   life insurance carriers                     & 526   & Q         & 36,719      \\ \hline
524114 & Direct health \& medical insurance   carriers        & 209   & Q         & 29,433      \\ \hline
52412  & Other   direct insurance carriers                    & 936   & Q         & 38,872      \\ \hline
524126 & Direct   property \& casualty insurance carriers     & 807   & Q         & 37,308      \\ \hline
524127 & Direct   title insurance carriers                    & 120   & Q         & (1000-2499) \\ \hline
524128 & All   other direct insurance carriers                & 9     & Q         & (100-249)   \\ \hline
52413  & Reinsurance   carriers                               & 63    & Q         & 2,709       \\ \hline
5242                & Agencies,   brokerages, \& other insurance related activities & 7,678                   & 9,646,288                  & 63,824                  \\ \hline
52421  & Insurance   agencies \& brokerages                   & 6,855 & 7,859,852 & 51,251      \\ \hline
52429  & Other   insurance related activities                 & 823   & 1,786,436 & 12,573      \\ \hline
524291 & Claims   adjusting                                   & 358   & 552,146   & 4,397       \\ \hline
524292              & Third   party administration of insurance \& pension funds    & 318                     & 1,017,610                  & 6,569                   \\ \hline
5242                & Agencies,   brokerages, \& other insurance related activities & 7,678                   & 9,646,288                  & 63,824                  \\ \hline
52421  & Insurance   agencies \& brokerages                   & 6,855 & 7,859,852 & 51,251      \\ \hline
52429  & Other   insurance related activities                 & 823   & 1,786,436 & 12,573      \\ \hline
524291 & Claims   adjusting                                   & 358   & 552,146   & 4,397       \\ \hline
524292              & Third   party administration of insurance \& pension funds    & 318                     & 1,017,610                  & 6,569                   \\ \hline
524298 & All   other insurance related activities             & 147   & 216,680   & 1,607       \\ \hline

\end{tabular}
\end{table}
\vfill

\raisebox{0pt}{\makebox[\linewidth]{\thepage}}

\end{landscape}
\clearpage


There was certainly at this moment, in Elizabeth's mind, a more gentle sensation towards the original than she had ever felt at the height of their acquaintance. The commendation bestowed on him by Mrs. Reynolds was of no trifling nature. What praise is more valuable than the praise of an intelligent servant? As a brother, a landlord, a master, she considered how many people's happiness were in his guardianship!—how much of pleasure or pain was it in his power to bestow!—how much of good or evil must be done by him! Every idea that had been brought forward by the housekeeper was favourable to his character, and as she stood before the canvas on which he was represented, and fixed his eyes upon herself, she thought of his regard with a deeper sentiment of gratitude than it had ever raised before; she remembered its warmth, and softened its impropriety of expression.


\subsection{Third-level subheading}

They were within twenty yards of each other, and so abrupt was his appearance, that it was impossible to avoid his sight. Their eyes instantly met, and the cheeks of both were overspread with the deepest blush. He absolutely started, and for a moment seemed immovable from surprise; but shortly recovering himself, advanced towards the party, and spoke to Elizabeth, if not in terms of perfect composure, at least of perfect civility.

She had instinctively turned away; but stopping on his approach, received his compliments with an embarrassment impossible to be overcome. Had his first appearance, or his resemblance to the picture they had just been examining, been insufficient to assure the other two that they now saw Mr. Darcy, the gardener's expression of surprise, on beholding his master, must immediately have told it. They stood a little aloof while he was talking to their niece, who, astonished and confused, scarcely dared lift her eyes to his face, and knew not what answer she returned to his civil inquiries after her family. Amazed at the alteration of his manner since they last parted, every sentence that he uttered was increasing her embarrassment; and every idea of the impropriety of her being found there recurring to her mind, the few minutes in which they continued were some of the most uncomfortable in her life. Nor did he seem much more at ease; when he spoke, his accent had none of its usual sedateness; and he repeated his inquiries as to the time of her having left Longbourn, and of her having stayed in Derbyshire, so often, and in so hurried a way, as plainly spoke the distraction of his thoughts.


% Place your additional chapters here using the \input{} command
\input{chapters/chapter_02}

\bibliographystyle{plain}
% \bibliographystyle{ieeetr.bst}
% \bibliographystyle{acm}
% \bibliographystyle{IEEEtran}

% Give this command the relative path to the .bib file.

\bibliography{references.bib}

% Appendices
% \appendix
% \chapter{Example Appendix 01}
\label{ch:appendix_01}

\blindtext
% \chapter{Example Appendix 02}
\label{ch:appendix_02}

\blindtext[2]

\end{document}
